% Part: second-order-logic
% Chapter: metatheory
% Section: undecidability-and-axiomatizability

\documentclass[../../../include/open-logic-section]{subfiles}

\begin{document}

\olfileid[hi]{sol}{met}{nax}

\olsection{द्वितीय-कोटि तर्कशास्त्र अभिगृहीतीकरणीय नहीं है}


\begin{thm}
\ollabel{thm:sol-undecidable}
द्वितीय-कोटि तर्कशास्त्र अनिर्णेय है।
\end{thm}

\begin{proof}
कोई प्रथम-कोटि !!{sentence}, प्रथम-कोटि तर्कशास्त्र में तभी और केवल तभी
वैध है जब वह द्वितीय-कोटि तर्कशास्त्र में वैध हो; और प्रथम-कोटि तर्कशास्त्र
अनिर्णेय है।
\end{proof}

\begin{thm}
\ollabel{cor:sol-not-axiomatizable}
द्वितीय-कोटि तर्कशास्त्र के लिए कोई ध्वनि और पूर्ण !!{derivation} पद्धति
नहीं है।
\end{thm}

\begin{proof}
मान लें $!A$ अंकगणित की भाषा का कोई !!a{sentence} है। $\Sat{N}{!A}$ तभी
और केवल तभी जब $\Th{PA^2} \Entails !A$। $!P$, ~$\Th{PA^2}$ के नौ
अभिगृहीतों का संयोजन हो। $\Th{PA^2} \Entails !A$ तभी और केवल तभी जब
$\Entails !P \lif !A$, अर्थात् $\Sat{M}{!P \lif !A$}। अब उस
!!{sentence} पर विचार करें
$\lforall[z][\lforall[u][\lforall[u'][\lforall[u''][\lforall[L][(!P' \lif !A')]]]]]$,
जो $\Obj{0}$ को $z$ से, $\prime$ को एक-स्थानी फलन-चर $u$ से, $+$ और
$\times$ को क्रमशः दो-स्थानी फलन-चरों $u'$ और $u''$ से तथा $<$ को
दो-स्थानी संबंध-चर~$L$ से बदलकर, फिर सार्वत्रिक परिमाणीकरण लगाकर प्राप्त
होता है। यह शुद्ध द्वितीय-कोटि तर्कशास्त्र का वैध वाक्य तभी और केवल तभी है
जब मूल वाक्य वैध था, तभी और केवल तभी जब $\Th{PA^2} \Entails !A$, और तभी
और केवल तभी जब $\Sat{N}{!A}$। अतः यदि द्वितीय-कोटि तर्कशास्त्र के लिए कोई
ध्वनि और पूर्ण प्रमाण-पद्धति होती, तो हम उसका उपयोग~$\Struct{N}$ में सत्य
!!{sentence} का संगणनीय परिगणन $f\colon \Nat \to \Sent[L_A]$ परिभाषित
करने के लिए कर सकते थे। यह फलन~$\Th{Q}$ में किसी प्रथम-कोटि
सूत्र~$!B_f(x, y)$ द्वारा निरूपणीय होता। तब !!{formula}
~$\lexists[x][!B_f(x, y)]$, ~$\Struct{N}$ के सत्य प्रथम-कोटि !!{sentence}
के समुच्चय को परिभाषित करता, जो टार्स्की के प्रमेय का खंडन है।
\end{proof}



\end{document}
